\section{Introduction}\label{sec:introduction}
A placement decision can be operationally attractive and still be
unacceptable. A node with spare CPU may lie outside a service's permitted
jurisdiction, lack encryption support, or fail an isolation requirement.
Once these restrictions are enforced, a second problem remains: a node
that looks lightly loaded now may be approaching a resource boundary.
Moving a service there can exchange one problematic destination for
another.

This paper studies the decision procedure at that boundary between policy
admission and resource management. We consider a controller that periodically
reconsiders existing microservice placements. Its inputs are service
requirements, node capabilities, and recent CPU and memory measurements.
Its output is a decision to retain the current source, select a different
node, or report that no candidate passes admission. The question is how to
use temporal information without allowing a favorable resource score to
override a mandatory policy condition.

Hard filtering before scoring is established practice. Kubernetes, for
example, separates filtering, scoring, reservation, and binding
\citep{kubernetes2026}. We do not claim that this separation is new.
Our focus is the interaction of source retention, a next-interval risk
screen, destination-selection history, and a separate headroom rule in
repeated placement decisions. This is narrower than solving general
edge--cloud orchestration: we do not optimize service-graph latency,
network traffic, or migration execution.

\cage uses four checks with different purposes. A Boolean policy mask
excludes prohibited nodes. A histogram gradient-boosting model uses recent
node history to estimate whether CPU or memory will cross a fixed boundary
in the next 30 seconds. Historical selection counts favor destinations
that have been used less often. Finally, a persistence projection with
empirical 95th-percentile residual margins checks headroom before a
destination is recorded. The source is retained whenever it passes the
policy, risk, and headroom conditions. Otherwise, the scheduler examines
a bounded shortlist and expands the search if necessary.

The distinction between prediction and admission matters. A probability
below the risk threshold is a model decision, not a safety guarantee.
Likewise, passing a headroom projection does not establish that a node can
absorb the incoming service: the projection does not add that service's
demand. These checks define the scheduler's decision contract. We evaluate
both whether it obeys that contract and whether selected nodes subsequently
look safer in the unmodified trace.

Our evaluation replays CPU-burst events from Alibaba's 2021 microservice
trace. Policy attributes are synthetic because the trace does not provide
them. On the natural test workload, \cage reduces migration decisions by
3.732 percentage points relative to a current-state learned baseline.
Against a policy-and-headroom baseline, it instead requests 0.920
percentage points more migrations, while the recorded next-state unsafe
rate falls from 0.155\% to 0.043\%. These are different trade-offs, not
evidence that one scheduler dominates all alternatives. The replay does
not execute migrations or modify future telemetry.

The paper makes three contributions:
\begin{itemize}[leftmargin=*]
\item A placement decision procedure that combines source retention,
hard policy admission, temporal risk screening, destination-use history,
and a separate empirical headroom check.
\item A replay comparison that distinguishes compliance with configured
checks from observed next-state behavior, including current-state and
component-removal controls.
\item An analysis of shortlist cost, policy-overlay sensitivity, and
the limits of conclusions drawn from unmodified trace replay.
\end{itemize}
